\documentclass[12pt]{article}
\usepackage{fontspec}
\usepackage[doublespacing]{setspace}
\usepackage[margin=1cm]{geometry}
\usepackage{fancyhdr} % <-- Add this line

% Set Arial as the main font
\setmainfont{Arial}
\setlength{\parindent}{0pt} % Remove paragraph indentation

% Add the following lines for page number positioning
\pagestyle{fancy}
\fancyhf{}

\rhead{\vspace{-.65cm}\thepage}
\begin{document}

Dear Admission Committee, 

Driven by a passion for AI-based autonomy, I conduct research on issue-oriented real-world challenges. My ambition is to enhance the safety of autonomous systems through certifiable, distributed robust autonomy for motion planning. I am particularly interested in enhancing the reliability of autonomous systems which require multi-agent planning under uncertainty, including extraterrestrial rovers and search and rescue robots. Several labs within NYU could support me in this research.
 
In collaboration with Dr. Feng, I endeavor to leverage my proficiency in LiDAR-based perception systems for the investigation and subsequent integration of LiDAR-derived models with multi-agent trajectory planners. This pursuit aims to culminate in the development of an advanced real-time 3D scene understanding system, fortified by certifiable perception algorithms explicitly engineered for GPS-denied environments that lack mapping infrastructure. The overarching objective is to augment this system with cutting-edge learning-based frameworks, endowing it with the ability to predict and adapt to environmental changes. This predictive capability is poised to empower robotic systems by enabling them to expeditiously formulate effective plans based on historical data and learned experiences. To further enhance semantic scene understanding, I advocate for the fusion of RGB-D cameras with LiDAR technology to improve performance and safety of autonomous vehicles, particularly in intricate urban settings.

With experience training visual navigation systems using various datasets, including the KITTI dataset, I bring substantial expertise to AI4CE where I will focus on enhancing LiDAR-based vSLAM techniques, aligning with the lab's overarching objectives of advancing learning-based robot perception, autonomous driving, and multi-agent systems to elevate computer vision for smart cities.


With Dr. Loianno, I explored state estimation and resilient controllers for safe autonomous drone flights. I believe my expertise can contribute to visual-inertial navigation projects. Actively involved in visual and visual-inertial navigation for unmanned vehicles, I concentrate on odometry and attitude estimation. My interest in robotic perception started with drone navigation projects, leading to the creation of the first open-source generalized inertial attitude estimation models.

At Machines in Motion, I aim to explore the foundational aspects of robotic movements, utilizing my expertise in reinforcement learning to advance legged locomotion. Legged robots hold significant potential for navigating challenging environments, proving valuable in planetary exploration, and search and rescue missions. I propose using Cooperative Reinforcement Inverse Learning (CIRL) with motion capture cameras to guide legged robots in mimicking human behavior. CIRL, by enabling collaborative learning from human demonstrations, addresses RL challenges like sample inefficiency and extensive trial and error. This approach is particularly beneficial in scenarios where human input shapes optimal behaviors. Immersed in research at the Machines in Motion lab, I anticipate conducting fundamental studies in robotics and applying this knowledge to real-world challenges.







Drawing on my varied experiences, I aim to contribute to Ph.D. research on robotic navigation and interaction. My expertise in state estimation and robotic perception will play a key role in developing precise semantic understanding for unmanned vehicles. I seek to apply these skills for safer operations, especially in critical situations requiring precision and adaptability. My goal is to make substantial contributions to the field of autonomous systems, ensuring their effectiveness and safety in real-world applications.


After earning my graduate degree from NYU Tandon, I plan to advance swarm robotics and autonomous systems for disaster response. Using cutting-edge algorithms, I seek to enhance coordination and communication among autonomous agents, facilitating quick deployment in complex environments. My goal is to contribute to the creation of a resilient autonomous swarm that can adapt dynamically to unforeseen challenges, improving disaster response capabilities and ensuring the safety of responders and the public. This aligns with NYU Tandon's emphasis on innovation and practical impact, offering a unique chance to translate theoretical concepts into real-world solutions.

My 5-year plan involves launching a startup dedicated to democratizing access to autonomous technologies. Inspired by my experiences as a first-generation student overcoming challenges, I aim to create accessible tools empowering individuals and organizations to leverage autonomous systems for diverse applications. This venture seeks to bridge the gap between academic research and practical implementation, making autonomous technologies a tangible reality across industries. Through this entrepreneurial endeavor, I aspire to contribute to the widespread adoption of autonomous systems in society, envisioning a future where innovation extends beyond labs and becomes an integral part of everyday life.

Driven by a passion for cultivating an inclusive and diverse STEM community, I, as a professor, am dedicated to empowering the next generation of scientists, regardless of their background. Being a first-generation student, I navigated unique adversities, grappling with limited academic resources and guidance. My commitment goes beyond technological contributions; it's about ensuring that those facing challenges akin to mine can enhance their lives and contribute meaningfully to our world. Through my fervor for teaching, I reinforce the dedication to creating an environment where everyone has the opportunity to thrive. This commitment is not merely a professional objective; it is a deeply personal mission.


In advance, I appreciate your time and consideration. 

Sincerely,

\end{document}